\section{Fate in the Ancient City}

\begin{quotex}
He who does not justly perform his appointed duty may appear as a violator of the whole order of the Universe.
\flright{\textsc{Pythagoras}}
\end{quotex}


In The Ancient City, \textbf{Fustel de Coulanges} begins in the middle of things, at the point when religion had already been established. To get to the beginning we have to turn to the poets. The Olympic gods with whom we are most familiar, were not even the first gods, but arose in time. Prior to the gods, the world was divided into three parts: Sky, Sea, and Earth. We can immediately see the relation to the Great Triad, as in the Taoist division into Heaven, Man, and Earth (the Sea represents the soul). Each of those parts was assigned to a god: Zeus, Poseidon, and Hades. This geographical apportionment of jurisdictions is carried over to the city. Each city has its god, as does the tribe, clan and family.

Since the god as founder of the city is its lawgiver, so the three brothers ruled their realms until Zeus claimed ultimate authority. Zeus then set down the law for all the gods and for men; the law was the result of the will of a god. Yet, there is a more primal principle. The Fates preceded the gods, who were even powerless against them. Even Zeus could not undo what the fates had destined for man.

\paragraph{Moira}


Before the gods, there was Moira, which we can understand as Destiny and is equivalent to the Tao or Dharma. Even the gods are limited by Destiny. Schopenhauer revived this idea with the notion of Will as the blind, purposeless, and determining force behind the phenomenal world. Nevertheless, that is only part of it. Moira is not just Fate, it is also Law in the sense of Right. Not only does it determine what must be, but also what ought to be, there is not only Destiny, but also Justice which involve the Will. These two meanings of law is dealt with in The Metaphysics of Law\footnote{\url{https://gornahoor.net/?p=447}}.

Therefore one's destiny is not merely an external limitation, but also the result of one's will: to fulfill one's duty or not. Thus, a man may be punished for his sins, but that is his own fault, not that of the Fates. The ancient Greeks did not have the notion of impossibility, but rather that of what lies beyond one's power, i.e., Privation\footnote{\url{https://gornahoor.net/?p=1815}}. This privation may be a self-limitation due more to a moral obligation than mere physical impossibility. This pagan understanding of the world is fundamentally moral. So the fundamental order of the world is not just the consequence of Destiny, but also of Will. Any action of men that disturbs this order, will have consequences. We can see the continuity of these notions with some Christian dogmas, properly understand, even though they are often denied by neopagans.

\begin{description}

\item[Hereditary Guilt]

\textbf{Aeschylus} wrote of the ``taints and troubles which, arising from some ancient wrath, existed in certain families.'' These were transmitted by blood from generation to generation. This notion served the Solidarity of the City; everyone ought to do, and must do, his duty. Otherwise, the family, or even the city, is affected unto later generations. (Some neopagans falsely deny that the ancient pagans could even experience guilt!)
\item[Moral corruption and nature]

Just as the fall of Adam altered the natural world, so did the pagans believe a similar thing. Hesiod wrote that when men do justice, their city flourishes and they are free from war and famine. Thebes was punished by the sin of Oedipus, according to the play by \textbf{Sophocles}.

\item[Vicarious sacrifice]

The ancient Greeks offered animal sacrifices to appease the gods. Some northern peoples even performed human sacrifices.
\end{description}

Part II will address the rise of philosophy and the twilight of the gods.\footnote{This part seems not to be ever written.}

Reference: \emph{From Religion to Philosophy} by \textsc{F. M. Cornford}.


\flrightit{Posted on 2011-08-09 by Cologero}

\begin{center}* * *\end{center}

\begin{footnotesize}
\begin{sffamily}

\texttt{Cologero on 2011-08-09 at 12:31 said:}

I apologize, Hoo, for giving the impression of an anti-Nordic sentiment, since it is far from out intention. If you focus on my own personal opinions, you will recall that I praised the Nordics for rebuilding Europe after the collapse of Rome. I praised the Normans for their all too brief stay in Sicily, where they maintained a fortress in my ancestral home of Salemi; I prefer to believe I retain some of their spirit, however little. You will find more.

These must be distinguished from translations, quotations, or historical references. I translated a paragraph on the Scandanavians from Evola, because of its historical interest. There is no doubt that Scandinavia is far along the leftist path today; what I hadn't realized is that even 70 years ago, that change was noted. I also quoted it to refute the overemphasis today on genetics, when a spiritual renewal is called for.

I mentioned sacrifice not to oppose it, on the contrary, the flow of blood is sometimes necessary however much the modern mind is displeased by the thought. I have mentioned animal sacrifice in connection with the Spartans to show their strength of will and commitment to their spiritual vision of the world. As for human sacrifice among the Germanics, this is quite different from, say, the brutality of the Aztecs, who sacrificed their enemies willy-nilly. No, the Germanics sacrificed the best among them, willingly and rarely. This is quite a different thing from what you describe and may be worth a thought some day.

So, Hoo, we oppose polemics in general, please read with a sense of detachment. The point of allowing comments is to invite clarifications and provide the opportunity for philosophical, historical, logical, or factual corrections. Nevertheless, we are forced to oppose modern movements that do not lead back to the primordial tradition, which is fundamentally a proto-Nordic tradition. Nietzsche is a superior psychologist, but metaphysics is deeper. We can discard the psychological and the historically contingent to penetrate to what may lie beneath. We never argue for or against any particular theological or historical position, although we may use them for illustrative purposes. At least we try, despite inevitable misunderstandings and the inability to express ourselves with sufficient clarity.

\hfill

\texttt{I on 2011-08-09 at 15:04 said:}

Interesting. What came to my mind while reading about guilt is that was it not Evola who distinquished between guilt and shame, and saw guilt as a Semitic emotion, while praising shame as an Aryan emotion? I think there is a sublime yet a very meaningful difference between these two emotions: for example, guilt arising from a (socio-)moral mis-behaviour and shame from a violation of one's true being.

Currently we are living largely in a civilization of (many times false) guilt, which is one typical phenomenon of internal judaification, whereas if we'd be living in a culture of shame, things would be quite differently both internally AND externally.

\hfill

\texttt{Cologero on 2011-08-09 at 18:03 said:}

We can agree about the misuse of guilt. If I understand your distinction, shame is something I bring on myself, whereas guilt is something imposed on me.

But to get back to Oedipus ... did he experience guilt or shame? The guilt is objective even if incest was not his intention. It could only be shame had he done it deliberately.


\hfill

\end{sffamily}
\end{footnotesize}

